\documentclass{situationist-preprint}

\title{Redemption as Redescription}
\msid{REL-22-b}
\msfamily{REL-22}
\msdiscipline{Philosophy of Religion}
\mstype{Consequence}
\mskeywords{redemption; defeat of evil; Chisholm; narrative sentences; significance; the past}

\begin{document}
\maketitle

\begin{abstract}
REL-22-a concludes that God cannot change the past but can change the significance of past events in light of later ones, and that past evils can thereby be defeated in Chisholm's sense without ceasing to be evils. We derive a reading of the language of redemption. Sentences such as \enquote{her suffering was not in vain} and \enquote{the betrayal was redeemed} ascribe significance properties, and they become true when later events make new descriptions of the past true. Three consequences follow: redemption never erases; it can be pending, its truth unsettled; and it must be distinguished from compensation, which corresponds to Chisholm's balancing off rather than to defeat.
\end{abstract}

\section{The reading}

\citet{Danto1965} called \emph{narrative sentences} those that describe an earlier event by reference to a later one, and observed that their truth can be settled only after the later event. REL-22-a identifies the properties such sentences ascribe as significance properties (Definition 2.2), and shows that an event can come to have one, as of a later time, without any change in the event or in anything earlier (Lemma 3.2 and Theorem 3.3). \citet{Chisholm1968} distinguished the balancing off of an evil, in which a whole containing it contains outweighing goods, from its defeat, in which the whole is better for containing the evil because of how the evil is related to the goods.

\begin{definition}\label{d:red}
A \emph{redemption sentence} is one that says of a past evil $e$ that it was redeemed, made good, or not in vain. It is \emph{true as of} $t$ if, as of $t$, $e$ is defeated in Chisholm's sense in a whole that includes later events.
\end{definition}

Being defeated in such a whole is a significance property, as the proof of Corollary 4.1 of REL-22-a observes. Redemption is thus redescription in a precise sense: what changes is which descriptions of the past are true, and it changes because later events occur, not because anyone retells the story.

\section{Three consequences}

\begin{proposition}\label{p:erase}
No redemption sentence entails that the evil did not occur or was not an evil.
\end{proposition}

\begin{proof}
By Corollary 4.1 of REL-22-a, a defeated evil does not cease to be an evil, and by Theorem 3.1 nothing can make it not have occurred. A redemption sentence ascribes only a relational property of the evil.
\end{proof}

The phrases of redemption that suggest erasure, that it is as if the wrong had never been done, are on this reading figures for something else: that the wrong no longer determines what the whole is. Read literally they are false, and they are not what redemption requires.

\begin{proposition}\label{p:pending}
A redemption sentence about $e$ may lack a settled truth value at a time after $e$, and acquire it later.
\end{proposition}

\begin{proof}
Whether $e$ is defeated depends on later events. Where these are future contingents at $t$, there are histories through $t$ on which $e$ is defeated and histories on which it is not (Assumption 2.3 of REL-22-a), and the sentence is not settled as of $t$ (Definition 2.4). Lemma 3.2 of REL-22-a shows how it can become settled later.
\end{proof}

Proposition~\ref{p:pending} gives the grammar of hope. \enquote{This will not have been in vain} is, at the time of the evil, neither known nor knowable; it is a commitment whose fulfillment lies in later events, some of them free acts of the sufferer and of others, as REL-22-a stresses.

\begin{proposition}\label{p:comp}
Compensation is not redemption.
\end{proposition}

\begin{proof}
Compensation adds goods to a whole containing the evil, so that the whole is good on balance. That is balancing off in Chisholm's sense. A whole in which an evil is balanced off would be at least as good without the evil, so the evil is not defeated, and by Definition~\ref{d:red} the redemption sentence is not made true.
\end{proof}

\begin{objection}
Redescription is a change in our descriptions, not in reality. A redemption that consists in new true sentences about the past is a redemption in the eye of the narrator.
\end{objection}

\begin{reply}
The new sentences are true because of later events, and a narrator who tells them before those events tells something false or unsettled. What is changed is not a description but the relational facts that make descriptions true. That those facts are relational does not make them less real than the battle that began a war, which was the beginning of the war only because the war followed.
\end{reply}

\section{Conclusion}

The language of redemption is literally true of the significance of past events. Redemption is defeat, not erasure and not compensation; it can be pending; and it is brought about by later events, God's and creatures', which make true of the past what was not true of it before. This is the sense in which God changes the significance of the past without changing the past.

\begin{thebibliography}{99}\small
\bibitem[Chisholm(1968)]{Chisholm1968} Chisholm, R.~M. (1968). The defeat of good and evil. \emph{Proceedings and Addresses of the American Philosophical Association} 42: 21--38.
\bibitem[Danto(1965)]{Danto1965} Danto, A.~C. (1965). \emph{Analytical Philosophy of History}. Cambridge University Press.
\end{thebibliography}

\end{document}
